\documentclass[10pt,twocolumn]{article}

\usepackage[T1]{fontenc}
\usepackage[utf8]{inputenc}
\usepackage{mathptmx}
\usepackage[margin=2cm]{geometry}
\usepackage{amsmath,amssymb}
\usepackage{booktabs}
\usepackage{titlesec}
\usepackage{fancyhdr}
\usepackage{enumitem}
\usepackage{siunitx}

\titleformat{\section}[block]{\centering\scshape\large}{\thesection.}{0.5em}{}
\titleformat{\subsection}[block]{\bfseries}{\thesubsection}{0.5em}{}
\setlength{\columnsep}{0.8cm}

\pagestyle{fancy}
\fancyhf{}
\fancyhead[R]{\small\itshape Formula 1 race and championship prediction with machine learning}
\renewcommand{\headrulewidth}{0pt}

\title{\bfseries Formula 1 race and championship prediction with machine learning}
\author{%
  Simone Bonasorta\\ Liceo Scientifico Nicola Pizi, Palmi
}
\date{}

\begin{document}
\twocolumn[%
  \begin{@twocolumnfalse}
    \maketitle
    \begin{center}\textbf{Abstract}\end{center}
    \noindent
    I made this project after the tennis one, to keep growing my knowledge of
    machine learning and to see if the same way of working would hold on a
    different sport. The idea is that the model works on thousands of race
    results collected from 1950 to 2026, the full history of the Formula~1
    World Championship. This time I built \emph{two} models, not one: a
    \textbf{pre-race} model that estimates, before the start, the probability
    that a driver finishes on the podium, and a \textbf{full-season} model that
    estimates, at any point of the year, the probability that a driver wins the
    Drivers' Championship. Both use only data that was known before the event.
    The strong part is again the features I extract and use. In this paper I show
    how I built the two models, how the pre-race one does against the qualifying
    grid, and how the probabilities relate to the ones used by the BookMaker, the
    AI that runs on betting sites. The pre-race model is above the qualifying
    grid, by a small margin on one season and by more when the seed is picked on
    the validation year. The season model needs careful calibration before its
    numbers mean anything. I keep both honest: the code and the trained models
    are public, so every number here can be reproduced.
    \vspace{0.6em}
  \end{@twocolumnfalse}
]

% ---------------------------------------------------------------
\begin{center}\scshape\large Indices\end{center}
\begin{enumerate}[leftmargin=1.4em,itemsep=2pt]
  \item[0.] \textbf{Abstract:} A short summary of the project and the two
            models I built.
  \item[1.] \textbf{Introduction:} Why I moved to Formula~1 after tennis, and
            how I started the project.
  \item[2.] \textbf{Related work:} The baseline I compare against and the role
            of the BookMaker.
  \item[3.] \textbf{Methodology:} The two models, architecture, dataset,
            preprocessing and features.
  \item[4.] \textbf{Results:} Accuracy, calibration and comparison with the
            grid baseline.
  \item[5.] \textbf{Discussion:} What works, the blocks I hit, and why some
            ideas did not help.
  \item[6.] \textbf{Conclusion:} Final remarks and what I take from the
            project.
  \item[7.] \textbf{References.}
\end{enumerate}

% ---------------------------------------------------------------
\section{Introduction}
I made this project after the tennis one, for the same amateur reasons and for
fun. The tennis model was my first real attempt at \emph{win} prediction: I
started, like most public tutorials, from a Random Forest, and I got a clear
improvement over the amateur baseline before hitting a plateau. For a first try
that was fair, and the plateau itself taught me something, because it was not
the data running out but my techniques.

So with Formula~1 I wanted to bring the knowledge I had gained and push past
that plateau on purpose. What carried the improvement, in the end, was not a
heavier model but three things on top of gradient boosted trees: richer
features built by hand, a time-decay weighting that makes recent seasons count
more, and a calibration step that turns the raw scores into honest
probabilities. I did try a heavier idea too, a temporal convolutional network
on a driver's recent races, but it did not make the final model; I keep that
story for the discussion, because the reason it failed is itself useful.

The way I started was practical. I am not a motorsport analyst, so first I used
an AI research assistant (Claude) to map out which public datasets exist for
Formula~1, how deep they go and what they contain. Once I had the sources, the
rest was on me. I loaded the data, looked at what was actually inside, and built
the features one by one by reasoning about what should matter for a race,
instead of throwing everything at the model and hoping.

My starting question was the same as in tennis: how high can I push the
accuracy using only free public data, and where is the real limit. But Formula~1
gave me a second question that tennis did not have. A season is a long story,
not a single match, so I also wanted to predict \emph{who wins the title},
updating the guess race after race.

As in the tennis paper, I want to be clear from the start. The output of both
models is a \emph{probability}, not a bet. I never placed real bets with it.
The right thing to compare these probabilities against is the grid for the
race, and the BookMaker for the title.

% ---------------------------------------------------------------
\section{Related Work}

\subsection{The baseline}
The reference I compare the pre-race model against is the \textbf{qualifying
grid}: the simple rule that the three drivers starting first are the three that
will finish on the podium. In Formula~1 this is a strong baseline, because track
position is hard to recover and pole-sitters win a lot. Any honest improvement
has to show as a clear gain over this rule, on races the model has not seen
during training.

My model does not classify a driver into fixed classes like ``winner / points /
no points''. It \emph{ranks} the drivers by their probability of finishing on
the podium, and is judged on how many of its top three guesses are correct. This
is the number that matters for the goal I set, so it is the one I report.

\subsection{The BookMaker}
The BookMaker is the same system I described in the tennis paper: the AI that
runs on betting sites and sets, for each outcome, a price that turns into a
probability once the margin is removed. For Formula~1 the natural market to
compare against is the outright market, where you bet, weeks in advance, on who
wins the championship. Given a decimal odd \(o\) for one driver, the implied
probability, before normalisation, is \(1/o\). Dividing each by the sum over all
drivers removes most of the margin and gives a probability that sums to one.

Unlike tennis, I could not find a free historical archive of Formula~1 odds. So
for this sport the BookMaker comparison is done live, on the 2026 season, by
writing down the championship odds of an Italian bookmaker (Sisal) and comparing
them with the season model week by week.

One clarification on framing. The odds are an efficient-market approximation, not
ground truth: they are not the "right" answer, just a strong, money-backed
estimate. So when I compare my probabilities to them I am doing a qualitative
calibration check (do my numbers live in the same region as a serious market?),
not claiming to beat the market. Beating a liquid market on free public data is
not the goal of this project.

% ---------------------------------------------------------------
\section{Methodology}

\subsection{Two models, one method}
I built two separate models that share the same pipeline but answer two
different questions.

\begin{itemize}[leftmargin=1.2em,itemsep=2pt]
  \item The \textbf{pre-race model} predicts, before a Grand Prix, the
        probability that a driver finishes in the top three. It is evaluated
        with \emph{precision@3}: among the three drivers with the highest
        predicted probability, how many actually reach the podium, averaged
        over the races.
  \item The \textbf{full-season model} predicts, after round~\(N\) of a season,
        the probability that a driver ends the year as champion. It is
        evaluated on how often its top pick is the real champion, and on how
        well its probabilities are calibrated.
\end{itemize}

Both work only on \textbf{pre-event} data: for the race, what is known before
the lights go out (grid, qualifying gap, recent form, car and driver history);
for the title, what is known after the last race already run (points, gap to the
leader, races left). I did not use any in-race telemetry or live timing.

\subsection{Architecture}
Both models share the same core: a gradient boosted tree ensemble, this time
\textbf{CatBoost} [2], with ordered boosting that behaves well on a dataset that
is not huge. They differ in how the raw scores are turned into probabilities.
The pre-race model puts an \textbf{isotonic regression} calibrator [4] on top,
fitted on the validation season, so the podium probabilities can be trusted and
compared to the grid. The season model uses a different step, described at the
end of this subsection.

Two design choices for the pre-race model came out of experiments, not from
theory:

\begin{itemize}[leftmargin=1.2em,itemsep=2pt]
  \item \textbf{Time-decay sample weights.} Formula~1 changes a lot from one
        season to the next: rules, cars, tyres. I weight each training race by
        \(\exp(-\text{age}/\tau)\) with \(\tau\) of about 1.5 seasons, so recent
        races count more. This was the single change that gave the biggest jump
        in the pre-race model.
  \item \textbf{Seed selection on validation.} The podium metric is a ranking
        over a small number of races, so it is noisy from one random seed to the
        next. I train several seeds and keep the one that is best on the
        \emph{validation} season, never on the test, so the choice is honest.
\end{itemize}

The season model does not use isotonic. Since only one driver can win the title,
a probability per driver is not enough: the numbers have to sum to one over the
drivers of that season. So I take the raw CatBoost scores and pass them through
a \textbf{softmax} [5] over the drivers of each race state, with a
\textbf{temperature} fitted on past seasons. Without this, the model was wildly
overconfident: it gave the leader a \(100\%\) chance after five races. With it
the same leader drops to a believable number.

\subsection{Data Collection}
The data comes from the Ergast / Jolpica Formula~1 database [1], free and
covering every race from 1950 to today. After downloading all the seasons I had
around \num{26000} driver-race results. From these I built two datasets:

\begin{itemize}[leftmargin=1.2em,itemsep=2pt]
  \item the \textbf{race dataset}, one row per driver per Grand Prix, used for
        the podium model;
  \item the \textbf{season dataset}, one row per driver per ``state of the
        championship after round \(N\)'', about \num{22000} rows, used for the
        title model.
\end{itemize}

The split is strictly in time, exactly as in tennis: old seasons for training,
one recent season for validation, and the most recent complete season (2025) as
the held-out test, with the part of 2026 already run kept aside for live
prediction. A random split would leak the future of a driver into the past, and
the accuracy would look good for the wrong reason.

\subsection{Data Preprocessing}
The key preprocessing step is the same one that mattered most in tennis: the
\textbf{state-before-event rule}. For every race the features are computed from
the state of the world \emph{before} the race, and the rolling statistics
(form, podium rate, reliability, championship points) are updated only after the
result is read. The code literally reads the state, writes the feature row, and
only then updates the counters. If you do these in the wrong order, the result
leaks into the features and the accuracy becomes fake.

For the season model I also had to fix a subtle bug: for the year still being
run (2026), the number of rounds is not the final one, so ``races left'' was
computed as zero. I set the expected calendar length for the current season so
that the remaining races are counted correctly.

\subsection{Pipeline}
The pre-race pipeline, from raw results to a calibrated podium probability, is:
\begin{verbatim}
download results (Ergast/Jolpica, 1950-2026)
sort races by season and round
for each race, in time order:
    read state -> grid, quali gap, driver and
                  constructor form, podium rate,
                  track history, lag, recovery
    write feature row
    update state with the real result   # only now
split by season: train / val / test
weight train races by exp(-age / tau)   # time-decay
train CatBoost on train, early stop on val
fit isotonic calibrator on val
evaluate precision@3, NDCG@3, Brier on test
\end{verbatim}
The order inside the loop is the state-before-event rule: a race's features are
written before its result updates the ratings and rolling counters. The season
model follows the same shape, with championship-state features and a
softmax-plus-temperature step instead of the isotonic calibrator.

\subsection{Features Used}
The race feature vector is made of a few groups, all reasoned out from how a
Grand Prix actually works.

\paragraph{Grid and qualifying.} The starting position is the strongest single
feature, because in Formula~1 track position is hard to recover. I also use the
qualifying gap to pole and the gap to the team mate, which separate a driver who
is genuinely fast from one who is just in a fast car.

\paragraph{Driver and car form.} Rolling averages over the last races of finish
position, points, podium rate and retirement rate, kept separately for the
driver and for the constructor. The car is half of the result in this sport, so
the constructor features carry real weight.

\paragraph{Track-specific history.} The average past result of the driver and of
the constructor at this circuit, and a split between street circuits and
permanent ones, because some cars and some drivers are clearly better on one
type than the other.

\paragraph{Lag and recovery.} The results of the last one, two and three races
as explicit features, and a ``recovery'' number, the average of grid minus
finishing position, that captures how much a driver tends to gain during a race.

For the season model the features are different in nature: current championship
points and position, gap to the leader (absolute and as a fraction), points per
race, podium count so far, momentum over the last three races, and a head-to-head
measure against the current leader.

\subsection{Code and data}
The data is the public Ergast / Jolpica database [1]. The race dataset has about
\num{26000} driver-race rows, the season dataset about \num{22000}. The feature
code, the time-decay weighting, the isotonic calibrator [4] and the
softmax-plus-temperature step [5] are in Python (CatBoost, scikit-learn). The
full pipeline, the per-race predictions of the 2025 test season, and the
trained models are published on GitHub
(\texttt{github.com/sussolinoss/sports-ml-predictions}, see the tagged release), so every
number in this paper can be reproduced from the same code and weights.

% ---------------------------------------------------------------
\section{Results}

\subsection{Evaluation metrics}
The pre-race model is judged with \textbf{precision@3}. For each race I take the
three drivers with the highest predicted probability and count how many of them
finish on the podium, then average over the races:
\begin{equation}
  \text{precision@3} = \frac{1}{R}\sum_{r=1}^{R}
       \frac{\#\{\text{top-3 picks on the podium in race } r\}}{3}.
\end{equation}
Since precision@3 only counts hits and ignores the order inside the top three, I
also report \textbf{NDCG@3}, a standard ranking metric that rewards putting the
actual podium drivers higher in the predicted order. For the quality of the
probabilities, in both models, I report the Brier score
[3], the mean squared error between the predicted probability and the result:
\begin{equation}
  \text{Brier} = \frac{1}{N}\sum_{i=1}^{N} (p_i - y_i)^2 ,
\end{equation}
and, for the season model, the expected calibration error (ECE), the average gap
between predicted confidence and observed frequency over probability bins. For
Brier and ECE, lower is better.

\subsection{Pre-race model}
On the held-out 2025 season (24 races) the calibrated CatBoost model reaches a
\textbf{precision@3 of 0.778}, against \textbf{0.750} for the qualifying grid
(the three drivers starting first). I also compare it to two simpler rules on
the same features: a logistic regression and a trivial rule that picks the three
drivers with the best podium rate over their last ten races, both at
\textbf{0.639}. The model is above the grid, and clearly above the linear and
the naive rule, which says the gain comes from non-linear structure in the
features and not from overfitting a strong baseline.

\begin{table}[t]
\centering
\caption{Pre-race model, test season 2025 (precision@3).}
\label{tab:race}
\begin{tabular}{lc}
\toprule
Rule & precision@3 \\
\midrule
Recent podium-rate (last 10) & 0.639 \\
Logistic regression (same features) & 0.639 \\
Qualifying grid              & 0.750 \\
My model, reproducible run (AUC 0.94, Brier 0.061, NDCG@3 0.826) & \textbf{0.778} \\
My model, best seed on validation & 0.833 \\
\bottomrule
\end{tabular}
\end{table}

The size of the test calls for caution, and I want to quantify it instead of
hand-waving. Bootstrapping the 24 test races (resampling races with replacement,
2000 times) gives a 95\% confidence interval on precision@3 of
\textbf{[0.708, 0.847]}. The grid baseline (0.750) sits inside that interval, so
the 0.778 is a \emph{suggestive, not statistically significant} improvement on a
single season. NDCG@3 is 0.826, which says the order inside the top three is
mostly right even when a pick misses. Picking the seed on the validation year
pushes the point estimate to 0.833, but that is a best-of selection, not the
typical run; the plain reproducible model is the honest headline. The era also
matters: with one or two dominant cars the podium is more predictable, so the
ceiling is higher than in a season where everyone wins.

\subsection{Ablation}
To see which features actually carry the model, I removed one feature group at a
time and retrained (Table~\ref{tab:ablation}). Only one group moves the number
clearly: removing grid and qualifying drops precision@3 by about 14 points. Every
other group sits within the bootstrap noise band, and removing some of them even
nudges the metric up on this one season, which is exactly what a wide confidence
interval predicts. The honest reading is that the qualifying signal is the
backbone, and the rest of the features help mostly through calibration and
robustness rather than through a clear, season-proof accuracy gain.

\begin{table}[t]
\centering
\caption{Ablation: precision@3 when one feature group is removed (test 2025,
full model 0.778).}
\label{tab:ablation}
\begin{tabular}{lcc}
\toprule
Removed group & precision@3 & $\Delta$ \\
\midrule
grid + qualifying & 0.639 & $-13.9$ \\
track history     & 0.778 & $\pm0$ \\
track-type        & 0.792 & $+1.4$ \\
lag (last 1--3)   & 0.833 & $+5.6$ \\
street + recovery & 0.833 & $+5.6$ \\
form              & 0.833 & $+5.6$ \\
\bottomrule
\end{tabular}
\end{table}

As a real held-out check, Table~\ref{tab:2026} puts the model's podium call next
to the actual podium for the first five races of the 2026 season. The model was
trained only on data up to 2025, so these are genuinely unseen races, and the
real podiums come from the results database, not from a simulation. The model
gets 9 of the 15 podium slots right, a precision@3 of 0.600. Lower than on the
2025 test, as expected on a small and more open sample.

\begin{table}[t]
\centering
\caption{Predicted vs.\ real podium, first five races of 2026 (held-out, model
trained up to 2025).}
\label{tab:2026}
\small
\begin{tabular}{llc}
\toprule
Race (predicted / real) & & Hits \\
\midrule
Australia & pred: russell, antonelli, verstappen & \\
          & real: russell, antonelli, leclerc    & 2/3 \\
China     & pred: antonelli, russell, hamilton   & \\
          & real: antonelli, russell, hamilton   & 3/3 \\
Japan     & pred: antonelli, russell, piastri    & \\
          & real: antonelli, piastri, leclerc    & 2/3 \\
Miami     & pred: antonelli, verstappen, leclerc & \\
          & real: antonelli, norris, piastri     & 1/3 \\
Canada    & pred: antonelli, russell, piastri    & \\
          & real: antonelli, hamilton, verstappen& 1/3 \\
\midrule
\multicolumn{2}{l}{\textbf{precision@3}} & \textbf{0.600} \\
\bottomrule
\end{tabular}
\end{table}

One thing I like about this table: the model never misses the podium completely.
It calls all three right once (China), two of three twice (Australia, Japan),
and at least one in every single race. On a sport as chaotic as Formula~1, with
crashes and safety cars, never going fully blind on a podium is a good sign.

\subsection{Full-season model}
The season model is where calibration mattered most. Before the softmax and
temperature step, the model gave the early leader a \(100\%\) title chance after
five races, which is obviously wrong. After calibration, on the 2025 test season
the probabilities are well behaved and the top pick is the real champion in about
23\% of the race states across the year. That number looks low until you remember
that early in a season the title is genuinely open, so a good model \emph{should}
be unsure. Table~\ref{tab:cal} gives the calibration numbers with and without the
temperature step.

\begin{table}[t]
\centering
\caption{Season model calibration on the 2025 test (lower is better).}
\label{tab:cal}
\begin{tabular}{lcc}
\toprule
Variant & Brier & ECE \\
\midrule
Softmax, no temperature  & 0.0502 & 0.0514 \\
Softmax + temperature (T = 1.77) & \textbf{0.0399} & \textbf{0.0338} \\
\bottomrule
\end{tabular}
\end{table}

Table~\ref{tab:champ} shows the model's \emph{forecast} for the 2026 championship
podium after round five: who it expects to finish in the top three of the final
standings. The season is not over, so this is a prediction, not a result.
Antonelli, the leader at that point, is a near lock for a podium, with Russell
and Hamilton close behind and Leclerc as the realistic outsider. The column
$P(\text{top-3})$ is the probability of that driver finishing top three in the
standings, and $P_{\text{norm}}$ is the same numbers scaled to sum to one across
drivers.

\begin{table}[t]
\centering
\caption{Forecast: 2026 championship podium after round five (season not yet
decided).}
\label{tab:champ}
\begin{tabular}{lcccc}
\toprule
Driver & Pos & Pts & $P(\text{top-3})$ & $P_{\text{norm}}$ \\
\midrule
Antonelli & 1 & 118 & 100.0\% & 30.9\% \\
Hamilton  & 3 & 61  & 81.8\%  & 25.3\% \\
Russell   & 2 & 67  & 81.8\%  & 25.3\% \\
Leclerc   & 4 & 58  & 60.0\%  & 18.5\% \\
rest      & --& --  & \(<\)0.5\% & \(<\)0.1\% \\
\bottomrule
\end{tabular}
\end{table}

For the title itself (not just the podium), the same calibration gives Antonelli
about a \textbf{79\%} chance after round five, against the \(100\%\) the raw
model wanted to give. On the same date the Italian bookmaker Sisal priced him
around \(1.55\) (about \(65\%\) after removing the margin), so the model is a
little more confident than the market but in the same region. This is a single
in-season snapshot, not a backtest: the proper comparison will be possible only
at the end of the season, once the odds have been recorded race by race.

% ---------------------------------------------------------------
\section{Discussion}

\subsection{What works}
Three things carried the pre-race model. The grid position, as expected, is the
backbone. The time-decay weighting, which makes recent seasons count more, was
the change that moved the number the most. And the careful, honest seed and
validation choice keeps the reported accuracy fair rather than lucky.

\subsection{The blocks I hit}
This project taught me as much from what \emph{did not} work as from what did.
A few blocks are worth writing down:

\begin{itemize}[leftmargin=1.2em,itemsep=2pt]
  \item \textbf{The TCN was dropped.} The temporal convolutional network on a
        driver's last races did learn a useful embedding, and before the
        time-decay weighting was added it improved the validation metric. After
        the decay weighting was in, the two carried the same signal (recent
        form) and the TCN no longer added anything: the podium accuracy on
        validation went slightly down, so it was removed from the final model.
        The embedding was real, but redundant once a simpler feature already
        encoded the same information.
  \item \textbf{Averaging many seeds hurt the ranking metric.} Averaging the
        probabilities of many models is supposed to be safer, but on
        precision@3, a metric that only cares about the order of the top three,
        the averaging smoothed the borderline cases and \emph{lowered} the
        accuracy. A single well-chosen seed did better.
  \item \textbf{CPU and GPU gave different trees.} Training the same model on
        GPU instead of CPU changed the splits enough to lose a couple of points
        on this small dataset, so I fixed the head model to CPU for
        reproducibility.
  \item \textbf{No free odds history.} For tennis I had years of closing odds.
        For Formula~1 the free archives do not exist, and the betting API I
        checked does not even list the sport. So the BookMaker comparison had to
        become a live, week-by-week capture instead of a clean backtest.
\end{itemize}

\subsection{Why the season model is the hard one}
The race model has a strong, almost mechanical signal in the grid. The season
model does not. Early in the year the title is open, and the only honest
prediction is an uncertain one. This is why the calibration step is not a detail
in that model, it is the core of it. A title prediction that says \(100\%\)
after five races is not bold, it is broken.

\subsection{What would break this model}
It is worth being explicit about where this model would fail, because the
features lean on assumptions that are only true in normal seasons:
\begin{itemize}[leftmargin=1.2em,itemsep=2pt]
  \item \textbf{A regulation reset.} The 2026 rule change (and the 2027 one) can
        reshuffle the order. The time-decay weighting helps it adapt faster, but
        the first races of a new era are the hardest, because the recent past is
        no longer a good guide.
  \item \textbf{Rain-heavy seasons.} The model has almost no wet-weather signal.
        A season with many wet races, where grid position matters less and chaos
        matters more, would hurt the precision@3.
  \item \textbf{A new dominant team appearing.} The constructor features are
        rolling, so a sudden rise (or fall) takes a few races to be reflected,
        and the model will lag during the transition.
  \item \textbf{A very open field.} The whole edge over the grid shrinks when no
        car dominates: the more random the podium, the closer any model gets to
        the baseline.
\end{itemize}

\subsection{Limitations and future work}
The first limitation is the size of the yearly test. With 24 races every metric
is noisy, and the only real cure is to wait for more seasons. The second is data
depth. I used results and standings, not the lap-by-lap telemetry that exists
from 2023 on, which could feed long-run pace and tyre features. A natural next
step is to add those, and to keep writing down the 2026 championship odds, so
that by the end of the season I finally have a clean book to compare the season
model against.

% ---------------------------------------------------------------
\section{Conclusion}
I built two Formula~1 prediction models on free public data. One is a pre-race
model that estimates the podium, the other a full-season model that estimates
the champion. The pre-race model reaches a precision@3 of about \textbf{0.778}
on the 2025 test season, above the \(0.750\) of the qualifying-grid baseline and
well above the \(0.639\) of a logistic regression on the same features; with the
seed picked on the validation year it reaches 0.833, but the reproducible number
is the honest one. The season model only becomes trustable after a
softmax-plus-temperature step that stops it from being overconfident.

As with tennis, the part I keep is the method more than the model. How to weight
data that ages, how to choose a seed without cheating, and how to calibrate a
multi-driver probability so it sums to one. The models are a small piece of
code. The discipline around them is what I carry to the next project.

% ---------------------------------------------------------------
\section{References}
\begin{enumerate}[leftmargin=1.6em,itemsep=1pt]
  \item Ergast Developer API and the Jolpica-F1 successor, \textit{Formula~1
        results database}, 1950--present.
  \item L.~Prokhorenkova, G.~Gusev, A.~Vorobev, A.~V.~Dorogush, A.~Gulin,
        ``CatBoost: unbiased boosting with categorical features,''
        \textit{NeurIPS}, 2018.
  \item G.~W.~Brier, ``Verification of forecasts expressed in terms of
        probability,'' \textit{Monthly Weather Review}, 1950.
  \item A.~Niculescu-Mizil, R.~Caruana, ``Predicting good probabilities with
        supervised learning,'' \textit{ICML}, 2005.
  \item C.~Guo, G.~Pleiss, Y.~Sun, K.~Q.~Weinberger, ``On calibration of modern
        neural networks,'' \textit{ICML}, 2017.
\end{enumerate}

\end{document}
